Frontier language-model agents resolve up to 99.4\% of tasks in a repaired software-engineering benchmark, and laboratories report agents working for hours on research. Yet in use, about 88\% of one laboratory's pull requests still need human steering or human-written code; the best of 16 models passes 61.8\% of workflows in web applications it builds from scratch; and an audit found fabrications in 70.4\% of 54 AI-authored conference submissions. Drawing on 288 papers and reports and 237 open-source repositories, this review joins two literatures that developed largely in parallel: AI for software engineering and automated science. We argue that verification, not capability, is now the binding constraint. Agents tend to succeed where an external signal they cannot edit decides success, and fail, often silently, where they grade their own work. We review the evidence from agent design to reliability and cost, then propose a framework, design principles and a research agenda.
